\chapter*{Abstract}
\addcontentsline{toc}{chapter}{Abstract}

Modern vehicles, and electric vehicles in particular, rely on the Controller Area Network (\gls{can}) bus to let their electronic control units exchange data. The \gls{can} protocol was designed for reliability on a closed bus, not for security: it has no built-in authentication or encryption. This makes it an open target for attacks such as message injection, denial-of-service flooding, and replay attacks.

This report presents \textbf{CANvisionNative}, a hybrid-architecture Intrusion Detection System (\gls{ids}) for vehicle \gls{can} bus networks. The system combines a Windows desktop client, built with \gls{wpf} and \gls{mvvm}, with a Python backend exposing a \gls{rest} \gls{api} built on FastAPI. The backend decodes raw \gls{can} frames using vehicle-specific \gls{dbc} files, extracts timing and payload features, and scores each frame with a fusion of four detection layers: a timing-anomaly layer, a payload-continuity layer, a per-vehicle Isolation Forest machine-learning layer, and a temporal-persistence layer. Detected anomalies are explained in natural language through an AI explanation pipeline that combines rule-based reasoning with a Large Language Model (\gls{llm}) provider chain.

The system supports three operating modes: replay of recorded logs (\gls{csv}, \gls{mf4}, candump, Vector \gls{can}), live simulation, and live hardware capture through an ELM327 or SocketCAN adapter. It was trained on more than seven million \gls{can} frames from eighteen vehicle profiles and validated on a real one-hour Kia EV6 recording containing 97,741 frames.

A significant part of this project consisted of a rigorous validation and correction phase. Four confirmed software defects were found and fixed: a vehicle-model selection bug that caused every replay to be scored with a generic model instead of the correct vehicle profile; a driving-state misclassification bug that labelled 99.93\% of a real driving session as ``parking''; a reason-attribution bug that let raw, unbounded diagnostic values dominate the displayed root-cause label instead of the actual fusion-score inputs; and a replay-parser inconsistency that caused some log formats to bypass the canonical detection pipeline. Each fix was proven with direct, reproducible evidence rather than by tuning thresholds. After the corrections, the alert count on the validation recording dropped from 15,945 to 11,971 (a 25\% reduction) while ML-score saturation dropped from 78\% to 0\%, with no artificial suppression of alerts.

This report documents the requirements, architecture, implementation, and validation of CANvisionNative, and shows that a desktop-class, explainable, multi-layer \gls{ids} can be built and rigorously verified for real vehicle \gls{can} bus data using commodity hardware and open datasets.

\vspace{0.5cm}
\noindent\textbf{Keywords:} Controller Area Network, Intrusion Detection System, Electric Vehicle, Machine Learning, Isolation Forest, Anomaly Detection, Large Language Model, Automotive Cybersecurity.

\cleardoublepage
